
\section{\trSerre{}的 $R_1+S_2$ 正规性判别法}

\trDepth{}的观念给出了一个有用的\trNormality{}判别法，它归功于\trSerre{}。

\textbf{26.3.1.}\ \emph{\trhDefinition{}。}\quad 设 $A$ 是一个\trNoether{}环，而 $k\in\mathbb{Z}^{\ge0}$。我们说 $A$ 具有\textbf{性质 $R_k$}（即 $A$ \textbf{在\trCodimension{} $\le k$ 处正则}（regular in codimension $\le k$），或者更不严格地说，$A$ \textbf{在\trCodimension{} $k$ 处正则}），若对每个\trCodimension{}至多为 $k$ 的\trPrimeIdeal{} $\mathfrak{p}\subset A$，$A_{\mathfrak{p}}$ 都是正则的。（在\trCodimension{} $1$ 处正则的重要性，我们已经屡次见到。）于是，一个\trNoether{}环是正则的，当且仅当它对一切 $k$ 都具有性质 $R_k$。\footnote{原书此处作「regular and only if」，当为「regular if and only if」之误（漏排 if），径改。——译者注}

我们说 $A$ 具有\textbf{性质 $S_k$}，若对每个\trPrimeIdeal{} $\mathfrak{p}\subset A$，\trLocalRing{} $A_{\mathfrak{p}}$ 的\trDepth{}都至少为 $\min(k,\dim A_{\mathfrak{p}})$——「诸局部环直到\trCodimension{} $k$ 为止都是 \trCM{}的，而在那之后其\trDepth{}都至少为 $k$」。于是，一个\trNoether{}环是 \trCM{}的，当且仅当它对一切 $k$ 都具有性质 $S_k$。

\textbf{26.3.A.（容易练习）}\quad 注意：一个\trNoether{}环 $A$ 总是平凡地具有性质 $S_0$。证明：一个\trNoether{}环 $A$ 具有性质 $R_0$，若 $\Spec A$ 是「一般既约的」（generically reduced）：即在其每个\trIrreducibleComponent{}的\trGenericPoint{}处既约。\footnote{原书此处作「has property \(R_0\) if Spec A is generically reduced」，据\mbox{练习~26.3.C} 疑当为「if and only if」，兹照录原书。——译者注}证明：一个\trNoether{}环 $A$ 具有性质 $S_1$，当且仅当 $\Spec A$ 没有\trEmbeddedPoint{}。（可能的提示：\mbox{练习~26.1.A}。）

\textbf{26.3.B.（练习：$S_2$ 的重要初等解读）}\quad 设 $A$ 是一个\trNoether{}\trIntegralDomain。证明：$A$ 具有性质 $S_2$，当且仅当对一切非零的 $g\in A$，$A/(g)$ 都没有\trEmbeddedPrime{}。

\textbf{26.3.C.（练习）}\quad 证明：一个\trNoether{}环 $A$ 没有幂零元（即 $\Spec A$ 既约），当且仅当它同时具有性质 $R_0$ 与 $S_1$。（提示：证明一个\trNoether{}概形是既约的，当且仅当它的每个\trIrreducibleComponent{}都一般既约，并且它没有\trEmbeddedPoint{}。）

把\mbox{练习~26.3.C} 中的下标各加一，便得到\trSerre{}的判别法。

\textbf{26.3.2. \trhTheorem{}（\trSerre{}正规性判别法）。}\ ——\ \emph{一个\trNoether{}环 $A$ 是\trNormal{}的，当且仅当它具有性质 $R_1$ 与 $S_2$。}

（回想 \S5.4.1：一个\trNoether{}环 $A$ 称为\trNormal{}的，若它在一切\trPrimeIdeal{} $\mathfrak{p}$ 处的\trLocalization{} $A_{\mathfrak{p}}$ 都是\trIntegrallyClosed\trIntegralDomain。）于是，\trNormality{}的失效可能有两个原因：它可能是 $R_1$ 的失效（这是我们已然知道的，来自\mbox{定理~13.5.8} 中 (a) 与 (g) 的等价性），也可能是 $S_2$ 更为微妙的失效。满足 $R_1$ 但不满足 $S_2$ 的簇的例子，见\mbox{例~26.3.4}（交于一点的两个平面）与\mbox{练习~26.3.F}（\trPinchedPlane）。

\textbf{26.3.D.（容易练习）}\quad 证明\trSerre{}正规性判别法（\mbox{定理~26.3.2}）。提示：仔细考察\mbox{命题~13.5.24} 与\mbox{命题~13.5.26}。

\textbf{26.3.3. 应用。}\quad 首先，\trSerre{}正规性判别法蕴含：\trCM{}概形是\trNormal{}的，当且仅当它们在\trCodimension{} $1$ 处正则。于是，要检验 $\mathbb{P}^n_k$ 中（或者更一般地，在任何光滑簇中）的超曲面（或更一般地，\trRegularEmbedding）的\trNormality，只须检验它们的\trSingularLocus{}的\trCodimension{}大于 $1$。（这些陈述之所以为真，你应当把其中的细节想清楚。）特别地，这给出了\mbox{练习~5.4.I}(b) 与 (c) 的一个新（更为复杂的）证明。

\textbf{26.3.E.（练习：概念演练）}\quad 证明：二维\trNormal\trVariety{}都是 \trCM{}的。

\textbf{26.3.4.}\ \emph{例：交于一点的两个平面。}\quad \mbox{练习~26.1.G} 中的簇 $X$（交于一点的两个平面）不是\trNormal{}的（为什么？），但它在原点之外处处正则。这就蕴含：$X$ 不具有性质 $S_2$（从而不是 \trCM{}的），而无须\mbox{练习~26.1.G} 的那些代数操作。

我们已然知道这个例子不是 \trCM{}的，但同样的想法可以表明：\trPinchedPlane{} $\Spec k[x^3,x^2,xy,y]$（出现在\mbox{练习~13.5.H} 中）不是 \trCM{}的。由于该环的「外在」描述，很难用别的方式来做到这一点。

\textbf{26.3.F.（练习：\trPinchedPlane{}不是 \trCM{}的）}\quad 设 $A$ 是 $k[x,y]$ 的子环 $k[x^3,x^2,xy,y]$。证明：$A$ 不是 \trCM{}的。提示：\mbox{练习~13.5.H} 已表明 $A$ 不是整闭的。

\textbf{26.3.5.}\ \emph{\trRegular\trLocalRing{}是整闭的。}\quad \trSerre{}正规性判别法可以用来证明：正则局部环是整闭的——而无须绕道困难的\mbox{事实~13.8.5}（\mbox{Auslander--Buchsbaum 定理}），即它们都是\trUFD。正则局部环都是 \trCM{}的（\S26.2.5），并且（由\mbox{练习~26.3.C}）因它们是整环（\mbox{定理~13.2.11}）而在\trCodimension{} $0$ 处正则，于是我们只须证明：正则局部环在\trCodimension{} $1$ 处正则。我们可以援引另一条困难的\mbox{事实~13.8.2}，即正则局部环的\trLocalization{}仍是正则的；但至少对于\trPerfectField{}上有限生成代数的\trLocalization，我们已经证明了这一点（见\mbox{定理~13.8.3} 与\mbox{练习~21.6.B}）。

\textbf{26.3.6.}\ \emph{\trhCaveat{}。}\quad 正如下面的练习所清楚表明的，条件 $S_2$ 是对\emph{一切}素理想施加的条件，而不仅仅是\trCodimension{}至多为 $2$ 的那些素理想。

\textbf{26.3.G.（练习）}\quad 举出一个满足 $R_1$、并且在一切\trCodimension{}至多为 $2$ 的点处都是 \trCM{}的簇的例子，它不是\trNormal{}的。

\textbf{26.3.7.}\ \extrahard\ \emph{射影正规性。}\quad 我们不会使用射影正规性这一概念，但仍要提它一笔，因为许多读者可能会与它打交道，而且把它介绍出来很容易。它是闭嵌入 $\pi\colon X\hookrightarrow\mathbb{P}^n_k$ 的一条性质。令 $\mathscr{L}=\pi^*\OO(1)$ 是给出这一闭嵌入的\trLineBundle（\S15.2），并且为方便起见，令 $R(\mathscr{L})_\bullet:=\bigoplus_{m\ge0}\Gamma(X,\mathscr{L}^{\otimes m})$ 为「$\mathscr{L}$ 的一切幂的截面」的\trGradedRing。于是我们说（定义 1）：$X\hookrightarrow\mathbb{P}^n_k$ 是\textbf{射影正规的}（projectively normal），若 $X$ 是正规的，并且对一切 $m\ge0$，$\mathrm{H}^0(\mathbb{P}^n_k,\OO(m))\to\mathrm{H}^0(X,\OO_X(m))$ 都是满射。换言之，对一切 $m\ge0$，$\mathscr{L}^{\otimes m}$ 的一切截面都「来自」多项式；不存在不「源出」于 $\mathrm{H}^0(X,\mathscr{L})$ 的「新」截面。等价地（定义 2）：经由 $\mathscr{L}=\pi^*\OO(1)$，$R(\mathscr{L})_\bullet$ 是多项式环 $k[x_0,\dots,x_n]$ 的一个商。你很快就能证明：这等价于如下陈述（定义 3）：（$X$ 是正规的，并且）对一切 $m\ge0$ 都有 $\mathrm{H}^1(\mathbb{P}^n_k,\mathscr{I}_X(m))=0$；也等价于如下陈述（定义 4）：$X$ 是正规的，$R(\mathscr{L})_\bullet$ 由次数 $1$ 生成，而且 $X\hookrightarrow\mathbb{P}^n_k$ 是由某个完全线性系给出的嵌入（假定它是非\trDegenerate{}的）。

\textbf{26.3.H.（练习）}\quad 证明：$\mathbb{P}^n_k$ 中的\trCompleteIntersection{}是射影正规的，当且仅当它们是正规的，当且仅当它们在\trCodimension{} $1$ 处正则。

证明 $X\hookrightarrow\mathbb{P}^n_k$ 是正规的，当且仅当其\trAffine\trCone{}是正规的，这并不太难；见 [Ha1, Ex.~II.5.14]。
